\subsection{Current available KGs and pipelines}

\newcolumntype{P}[1]{>{\raggedright\arraybackslash}p{#1}}

\begin{table*}[!t]
\centering
\caption{Scholarly knowledge graphs compared on infrastructure dimensions. Coverage figures are snapshot-specific and use units native to each system (triples for RDF graphs, works/papers for property graphs, citation links for OpenCitations). Licences reflect data, not code.}
\label{tab:kg-infra}
\footnotesize
\setlength{\tabcolsep}{4pt}
\begin{tabular}{P{2.4cm}P{2.8cm}P{2.4cm}P{2.5cm}P{2.2cm}P{1.4cm}}
\hline
\textbf{System} & \textbf{Coverage (snapshot)} & \textbf{Schema} & \textbf{Query interface} & \textbf{KGE released?} & \textbf{Licence} \\
\hline
MAKG~\cite{farber2019makg,farber2022makgenhanced} & 8\,B triples (2019); 239\,M papers (2022) & RDF, SPAR + Dublin Core + PRISM + FOAF & SPARQL, Zenodo dump & Yes (TransE\slash TransR\slash DistMult\slash ComplEx\slash RESCAL) & ODC-BY \\
OpenAlex~\cite{priem2022openalex} & 209\,M works (2022) & JSON, bespoke; Concepts hierarchy & REST, S3 dump & No official & CC0 \\
SemOpenAlex~\cite{farber2023semopenalex} & 26\,B triples, 249\,M works (2023) & RDF, SPAR + SKOS + custom \texttt{soa:} & SPARQL, AWS S3 TriG & Yes (TransE\slash DistMult\slash ComplEx\slash GraphSAGE\slash GAT) & CC0 \\
S2AG~\cite{kinney2023semanticscholar,ammar2018construction} & 225\,M papers, 2.8\,B edges (2025) & Property graph, JSON & REST Graph API & SPECTER\slash SPECTER2 (doc embeddings) & ODC-BY \\
OAG / OAG-Bench~\cite{zhang2019oag,zhang2024oagbench} & 0.7\,B entities, 2\,B relations (v1, 2019) & JSON-line property graph & Bulk download & OAG-BERT (LM, not KGE) & Research-only \\
AceKG~\cite{wang2018acekg} & 3.13\,B triples, 62\,M papers (2018, static) & RDF, bespoke vocabulary & Turtle dump only & Only as baselines & Free access \\
ORKG~\cite{jaradeh2019orkg} & Curated; small scale & Contribution nodes, Neo4j + RDF export & REST, SPARQL & None published & CC BY-SA \\
OpenCitations~\cite{peroni2020opencitations,massari2024ocmeta,heibi2024ocindex} & 2\,B+ citation links, 98\,M+ works (2024) & RDF, SPAR + PROV-O & SPARQL, REST, Figshare & None official & CC0 \\
\hline
\end{tabular}
\end{table*}

The existing landscape of scholarly retrieval systems stratifies along three
axes. Retrieval is either text-centric or graph-native, queries are either
natural language or structured, and evaluation maturity varies widely across
systems. No current system sits at
the intersection a KG-native hybrid pipeline requires. The three subsubsections
below substantiate this from three angles, namely the data layer
(Section~\ref{subsec:kg-infrastructure}), the traditional search engines built
on it (Section~\ref{subsec:traditional-search}), and the LLM-augmented
assistants that consume both (Section~\ref{subsec:llm-agentic}).

\subsubsection{Scholarly knowledge graphs as data infrastructure}
\label{subsec:kg-infrastructure}

OpenAlex is the canonical ingestion target for any new scholarly KG, but its
non-RDF design forces stack-level choices that shape everything downstream.
Microsoft Academic Graph's retirement at the end of 2021 left
OpenAlex~\cite{priem2022openalex} as the de facto open successor at scale, and
its canonical interface is a JSON REST API with no SPARQL endpoint and no
reuse of the SPAR ontologies.  The remaining systems matter to this thesis as
design precedents rather than candidates.
SemOpenAlex~\cite{farber2023semopenalex} re-lifts OpenAlex into RDF to close
exactly that gap, inheriting its anchor from the MAKG
lineage~\cite{farber2019makg,farber2022makgenhanced}.  A property-graph
cluster comprises S2AG~\cite{kinney2023semanticscholar} (which extends Ammar
et al.'s literature graph~\cite{ammar2018construction}),
OAG~\cite{zhang2019oag,zhang2024oagbench}, and AceKG~\cite{wang2018acekg}.
ORKG~\cite{jaradeh2019orkg} offers curated semantic depth, and
OpenCitations~\cite{peroni2020opencitations,heibi2019coci,massari2024ocmeta,heibi2024ocindex}
models citations as first-class resources.  Table~\ref{tab:kg-infra} carries
the per-system detail on coverage, schema, query interface, embeddings, and
licence.

None of these systems defines a retrieval task, which is why this thesis
treats them as substrate rather than baseline.
OAG-Bench~\cite{zhang2024oagbench} comes closest but evaluates graph-mining
tasks (name disambiguation, paper source tracing) rather than end-user search.
Reported counts also drift substantially across snapshots. S2AG and OpenAlex
have moved since their primary publications, and ORKG publishes no content
counts in the peer-reviewed paper, so version-specific citation is the only
defensible practice. The data layer therefore supplies a substrate, not a
baseline, and this thesis ingests OpenAlex into a custom typed graph rather
than consuming any of these systems as-is.

\subsubsection{Traditional scholarly search engines}
\label{subsec:traditional-search}

Today's deployed search engines practise text-centric and citation-graph
retrieval separately, and no widely-used system seriously combines them under
a single ranking regime. Text-centric engines dominate daily use. Google
Scholar is the largest by record count, with roughly 389\,M
records~\cite{gusenbauer2019gs} and the highest citation coverage in independent
comparison~\cite{martinmartin2021coverage}, but its 1{,}000-result cap, weak
Boolean support and opaque ranking rule it out as a systematic-review
engine~\cite{gusenbauerhaddaway2020}. Among major open-access alternatives at
100\,M+-document scale, Dimensions~\cite{herzog2020dimensions},
Lens~\cite{jefferson2018lens} and OpenAlex-native
search~\cite{priem2022openalex} offer Boolean keyword search with faceted
concept filters, while CORE~\cite{knoth2012core,knoth2023core} prioritises
bulk OA full-text aggregation and underpins LLM pipelines such as
OpenScholar~\cite{asai2024openscholar}. None of these engines exposes the
citation graph as a queryable structure.

The closest deployed approaches to a text-plus-citation hybrid are an
embedding layer and a citation-intent classifier, and neither amounts to the
queryable structure this thesis requires.  Semantic Scholar layers
SPECTER~\cite{cohan2020specter} and SPECTER2~\cite{singh2023scirepeval}
citation-informed transformer embeddings over its literature
graph~\cite{ammar2018construction,kinney2023semanticscholar}, with
SciNCL~\cite{ostendorff2022scincl} a stronger contrastive successor on
SciDocs, so the citation graph supervises training rather than answering
queries.  Scite~\cite{nicholson2021scite} classifies citation statements as
supporting, contrasting or mentioning via a SciBERT classifier, and a
third-party audit~\cite{bakker2023evaluating} reports F-measures between 0
and 0.58 against the vendor's claimed 80\%+ precision, which illustrates the
broader reliability problem with vendor-reported numbers in this cluster.

Graph-native discovery tools treat seed-paper exploration as a UI problem
rather than a retrieval task. Connected
Papers~\cite{smolyansky2020announcing} computes co-citation and
bibliographic-coupling similarity over a $\sim$50k-candidate Semantic Scholar
subset and displays a force-directed graph. Research
Rabbit~\cite{researchrabbit} iteratively expands collections over Semantic
Scholar and OpenAlex. Inciteful~\cite{inciteful_docs}, since 2023 primarily
OpenAlex-backed, builds depth-2 local citation graphs and runs PageRank and
link prediction. None of the three has a peer-reviewed benchmark. Evaluations
are confined to library-guide and scientometric-blog
commentary~\cite{tay2021literaturemapping}. The absence of formal evaluation
is itself a design-space observation. The scientometric community has not yet
settled on a benchmark for graph-native scholarly discovery.

Provenance concentrates risk across this whole layer. Crossref, PubMed and
the legacy MAG feed OpenAlex~\cite{priem2022openalex} and Semantic
Scholar~\cite{kinney2023semanticscholar}, which in turn feed Connected
Papers, Research Rabbit, Inciteful, and most of the LLM-augmented systems in
Section~\ref{subsec:llm-agentic}. Citation coverage inherited from Crossref
and MAG is incomplete~\cite{visser2021largescale}, and OpenAlex's and
Semantic Scholar's concept and topic taxonomies both trace back to MAG's
field-of-study classifier. Any new system built on this stack inherits the
same gaps.

\subsubsection{LLM-augmented and agentic scholarly search systems}
\label{subsec:llm-agentic}

Layered on top of these retrieval substrates, a different category of system
has emerged since late 2023 that consumes them rather than competing with
them. These LLM-augmented assistants differ from the retrieval engines above
in output contract rather than architecture. They produce long-form
citation-backed answers via retrieval-augmented generation rather than ranked
paper lists. Representative
open systems include OpenScholar~\cite{asai2024openscholar},
PaperQA2~\cite{skarlinski2024paperqa2} and Ai2 Scholar
QA~\cite{singh2025ai2scholarqa}, each combining dense or hybrid retrieval
over scholarly corpora with multi-stage LLM synthesis, with ORKG
Ask~\cite{oelen2024orkgask,oelen2026orkgask} the closest neuro-symbolic
relative. Commercial SaaS equivalents (Elicit, Consensus, Undermind) sit
alongside these but lack peer-reviewed primary papers. Evaluation here is
dominated by long-form QA benchmarks such as LitQA2 within
LAB-Bench~\cite{laurent2024labbench} and ScholarQABench, measuring answer
correctness and citation faithfulness rather than retrieval quality. These
systems are included for landscape completeness rather than as direct
comparators, and the benchmarks they rely on, alongside those used by the
retrieval systems above, are surveyed in the next section.

\FloatBarrier